% !TEX root = master.tex
% !TEX program = lualatex
\lecture{4}{12-03-2025}{Laurent series and Taylor series}{}

\subsection{Extended Cauchy Integral Formula}
Under the same assumptions, for every integer $n \ge 0$ we have

\begin{corollary}
\[
f^{(n)}(z_0)
=
\frac{n!}{2\pi i}
\int_\Gamma
\frac{f(z)}{(z-z_0)^{n+1}}\,dz .
\]
\end{corollary}
Equivalently,

\[
\int_\Gamma
\frac{f(z)}{(z-z_0)^{n+1}}\,dz
=
\frac{2\pi i}{n!}f^{(n)}(z_0).
\]

\textbf{Remark.}  
This formula shows that holomorphic functions are infinitely differentiable.

\paragraph{Examples}
Let $f(z) = \frac{\cos z}{z}$
and let $\Gamma$ be a positively oriented simple closed curve.

We distinguish three cases:

\begin{itemize}

\item If $0$ lies on $\Gamma$, the integral is not defined.

\item If $0$ is outside $\Gamma$, then

\[
\int_\Gamma \frac{\cos z}{z}\,dz = 0
\]

by Cauchy's theorem.

\item If $0$ is inside $\Gamma$, then using Cauchy's integral formula

\[
\int_\Gamma \frac{\cos z}{z}\,dz
=
2\pi i \cos(0)
=
2\pi i .
\]

\end{itemize}

Let $\Gamma = \{z \in \mathbb{C} : |z-2| = 1\}$ oriented counterclockwise.  
Compute

\[
\int_\Gamma \frac{e^z}{(z-2)^3}\,dz .
\]

Using the extended Cauchy formula with $f(z)=e^z$, $z_0=2$ and $n=2$,

\[
\int_\Gamma \frac{e^z}{(z-2)^3}\,dz
=
\frac{2\pi i}{2!} f^{(2)}(2).
\]

Since

\[
f^{(2)}(z_0)=e^2,
\]

we obtain

\[
\int_\Gamma \frac{e^z}{(z-2)^3}\,dz
=
\frac{2\pi i}{2}e^2
=
\pi i e^2.
\]

\chapter{Laurent series and Taylor series}\label{chap:laurent_series_and_taylor_series} % (fold)

\section{Taylor Series}

Let $D \subset C$ open, $z_0 \in D$, and $R > 0$ such, that the disk $\{z \in C : |z - z_0| < r \}$ is contained in $D$. If $f: D \to C$ is holomorphic, it has a Taylor series expansion for every $z \in C$ with $|z - z_0| < R$
\[
\sum_{n=0}^{\infty}
\frac{f^{(n)}(x_0)}{n!}(x-x_0)^n.
\]
As Taylor series around a point is unique, one can compute Taylor series using known series expansions.

The largest possible $R$ is called \textbf{radius of convergence} (It's the distance between $z_0$ and the nearest singularity (not $z_0$).

\paragraph{Example}
For the exponential function,
\[
e^x
=
\sum_{n=0}^{\infty}
\frac{x^n}{n!}.
\]

This series converges for all $x \in \mathbb{R}$. \\
\textbf{Remark:} Functions $f$ s.t the Taylor series converges to $f$ in a neighborood of $x_0$ are called \textbf{analytic}.

\subsection{Taylor coefficients}
\begin{definition}[Taylor Coefficients]
  The Taylor coefficients of a function $f(z)$ around a point $z_0$ are defined as: \[
    c_n = \frac{f^{(n)}(z_0)}{n!}
  \]
  where $f^{(n)}(z_0)$ is the $n$-th derivative of $f$ evaluated at $z_0$.

  This comes from assuming $f$ has a power series representation: \[
    f(z) = \sum_{n=0}^{\infty}c_n(z-z_0)^n  
  \]
\end{definition}

\textbf{Remark:}

$(1)$ All holo functions are analytic. 

$(2)$ In practice: $R$ is the distance of $z_0$ from the $nearest$ point where $f$ is not holo or not defined. 

$(3)$ We can connect the Taylor series to Cauchy's integral formula \[
  f(z) = \sum_{n = 0}^{\infty} \frac{f^{\left(n\right)}(z_0)}{n!}(z-z_0)^n = \sum_{n=0}^{\infty}c_n(z-z_0)^n 
\]
where \[
  c_n = \frac{f^{\left(n\right)}(z_0)}{n!} = \frac{1}{2 \pi i}\int_{\gamma} \frac{f(\zeta)}{\zeta - z_0} d\zeta
\]
for every simple, close curve $\gamma$ aroung $z_0$ with int $\gamma \subset D$ and counterclockwise oriented.
\section{Laurent series}

\begin{theorem}[Laurent series]
 Let $D \subset \mathbb{C}$ be open and let $z_0 \in D$.  
Assume $f : D \setminus \{z_0\} \to \mathbb{C}$ is holomorphic.

Let $R>0$ such that $B_R(z_0) \subset D$.

Then for all $z \in B_R(z_0)\setminus\{z_0\}$

\[
f(z)=\sum_{n=-\infty}^{\infty} c_n (z-z_0)^n
\]

where for every simple closed counterclockwise oriented curve $\gamma$
such that $\mathrm{int}(\gamma)\subset D$ and $z_0$ is enclosed by $\gamma$,

\[
c_n=\frac{1}{2\pi i}\int_\gamma
\frac{f(\zeta)}{(\zeta-z_0)^{n+1}}\,d\zeta .
\]

\end{theorem}
\paragraph{Remark}

\begin{itemize}

\item The Laurent series is like a Taylor series but it may include
terms with polynomial singularities (negative powers).

\item As for Taylor series, the Laurent series representation is unique.
Hence if we determine the coefficients $c_n$ by any method
(not necessarily by path integrals), we have determined the Laurent series.

\item We do not assume that $f$ is holomorphic at $z_0$.
It might be holomorphic there or not; no hypothesis is made.

\end{itemize}

\begin{definition}[]
 Any Laurent series can be split into a \textbf{singular part} and a
\textbf{regular part}:

\[
\sum_{n=-\infty}^{\infty} c_n (z-z_0)^n
=
\underbrace{\sum_{n=1}^{\infty} c_{-n}(z-z_0)^{-n}}_{\text{singular part}}
+
\underbrace{\sum_{n=0}^{\infty} c_n (z-z_0)^n}_{\text{regular part}} .
\]

 
\end{definition}
In other words, a Taylor series is a Laurent series whose singular
part vanishes.

\paragraph{Example}

If $f : D \to \mathbb{C}$ is holomorphic and $z_0 \in D$,
then its Taylor expansion

\[
f(z)=\sum_{n=0}^{\infty} c_n (z-z_0)^n
\]

is a Laurent series whose singular part is zero since for all $n\leq -1$\[
  \frac{f(\zeta)}{(\zeta-z_0)^{n+1}} \text{is holo. near} z_0
\]
so by Cauchy's integral formula $\int_{\gamma} \frac{f(\zeta)}{(\zeta-z_0)^{n+1}} d\zeta = 0$
\paragraph{Example}

For the following function $f: \mathbb{C}\setminus{0} \to \mathbb{C}$ let find the Laurent series around $z_0 = 0$


1) $f(z)=\frac{1}{z}$, this is already a Laurent serieswith
\[
c_{-1}=1, \qquad c_n=0 \ \text{otherwise}.
\]

2) $f(z) = \frac{5z + 1}{z^2} = \frac{1}{z^2} + \frac{5}{z} = \sum_{n=-\infty}^{\infty}c_nz^n$ where \[
  c_n = 
  \begin{cases}
    1 \quad \text{if }n = -2 \\
    5 \quad \text{if }n = -1 \\
    0 \quad otherwise
  \end{cases}
\]
